\songtitle{El mejor delito}

\tag*{1996}
\begin{notes}
\end{notes}
\begin{style}
Traditional Colombian Caribbean Puya, A fast-paced, driving 6/8 rhythm executed by traditional percussion: the tambor alegre provides syncopated, improvisational improvisations, the llamador keeps a strict offbeat counter-tempo, and the tambora provides deep, resonant downbeats along with the metallic scraping of the maraca or guache, The melodic core consists entirely of a duet between the gaita hembra, which leads with rapid, piercing, and ornamental melodic runs, and the gaita macho, playing lower-registered harmonic support and marking the beat with its single finger hole, A solo male vocal delivers the décima espinela with a rustic, powerful, and syllabic delivery, phrasing the ten lines tightly to match the rapid tempo, The tempo is a brisk 140 BPM in 6/8 time, The mix is raw and acoustic, with an organic, open-air room reverb that emulates a live, rural performance
\end{style}
\begin{lyrics}
\songpart{Verse 1}
El gris del cielo nublado\\
alivia el encierro frío,\\
mas crece un vacío impío\\
sobre el banco del cansado.\\
El verde en el diagramado\\
miente su fiel trayectoria,\\
y afuera la vida en gloria\\
invita a salir, a ser,\\
pero él decidió saber\\
que hay algo en esta memoria.

\songpart{Chorus}
No puede llorar del todo,\\
que la risa le hace guerra,\\
y en la muerte de la tierra\\
brota un brote de algún modo.

Odia no poder odiar\\
sin que el bien lo contrabalance,\\
que en cada grieta un romance\\
le impide el bien derrumbar.

\songpart{Verse 2}
Se toma un respiro inerte,\\
la mente en blanco a propósito,\\
mas vuelve el mismo compromiso\\
que el alma vuelve más fuerte.\\
Y en ese espacio advierte\\
la voz que quisiera oír:
"¿Cómo 'tash?"—debe existir—\\
y el "vale" de Andrea, mas ella\\
en la finca de aquella estrella,\\
y él sin teléfono, a sufrir.

\songpart{Chorus}
No puede llorar del todo,\\
que la risa le hace guerra,\\
y en la muerte de la tierra\\
brota un brote de algún modo.

Odia no poder odiar\\
sin que el bien lo contrabalance,\\
que en cada grieta un romance\\
le impide el bien derrumbar.

\songpart{Verse 3}
Odia el instante que clara\\
la verdad le hace mirar:\\
no puede del todo odiar\\
la vida que lo ampara.\\
Su carrera, aunque amarga\\
en domingos de circuito,\\
le ha dado el mejor delito:\\
amigos, risas, café,\\
y aunque el gris lo tenga en pie,\\
sabe que es su propio rito.

\songpart{Chorus}
No puede llorar del todo,\\
que la risa le hace guerra,\\
y en la muerte de la tierra\\
brota un brote de algún modo.

Y así ovidó odiar\\
en un temporal balance,\\
y en cada grieta un romance\\
le impide el bien derrumbar.
\end{lyrics}

